% GENERATED FILE. Edit canonical lessons, then run npm run generate:books.
% canonical-source-hash: fnv1a64:f09f2d147aa446de
% canonical-lessons: JA-C230-shachou, JA-C230-joushi, JA-C230-senpai, JA-C230-kouhai, JA-C230-ganbaru

\chapter{Bosses, Seniors and Juniors}
\label{ch:ja-bosses-seniors-and-juniors}
\hlchaptermodality{230}

\begin{chapteropening}
\textbf{By the end of this chapter:} \emph{I can talk about the people above and below me at work.}

Complete the last lesson of chapter 230: I can talk about the people above and below me at work.
\end{chapteropening}

\section[\texorpdfstring{\ja{しゃちょう} (shachō)}{しゃちょう (shachō)}]{\ja{しゃちょう} (shachō) — a company president}

\label{lesson:JA-C230-shachou}



\begin{quote}
Before the new one: say the Japanese for to be employed, then the Japanese for a colleague.
\end{quote}

\subsection*{You'll want to know: \ja{しゃちょう}}
\textbf{\ja{しゃちょう}} — \emph{shachō} — \textquotedblleft{}a company president\textquotedblright{}.

\begin{quote}
\textbf{\ja{しゃちょうは} \ja{いつも} \ja{かいぎで} \ja{はなします。}} — \emph{shachō wa itsumo kaigi de hanashimasu} — \textquotedblleft{}The company president always speaks at meetings.\textquotedblright{}
\end{quote}

\begin{cousinweb}
\textbf{\ja{しゃちょう}} (\emph{shachō}) is Sino-Japanese (kango): two Chinese characters read with their on'yomi, readings borrowed from Chinese. They are \emph{sha}, \textquotedblleft{}company, association\textquotedblright{} (Middle Chinese \emph{dzyaeX}) and \emph{chō}, \textquotedblleft{}leader, senior\textquotedblright{} (\emph{trjangX}) (KANJIDIC2; Middle Chinese in Baxter's transcription).
\end{cousinweb}

\begin{grammarlens}[title={What you've built}]
One more word for everyday life at A2.
\end{grammarlens}

\subsection*{Your turn}
\begin{itemize}
\raggedright
  \item \emph{Say it:} \emph{shachō}
  \item \emph{Say it:} \emph{shachō}, once more
  \item \emph{Say it:} \emph{shachō wa itsumo kaigi de hanashimasu} — The company president always speaks at meetings.
  \item \emph{Recall:} say the Japanese for to be employed, then the Japanese for a colleague, then say \emph{shachō} again
\end{itemize}

\begin{tcolorbox}[breakable,colback=teal!4,colframe=teal!35!black,title={Before you move on}]
What does \ja{しゃちょう} mean? (\textquotedblleft{}A company president\textquotedblright{}.) Say it once more.
\end{tcolorbox}
\section[\texorpdfstring{\ja{じょうし} (jōshi)}{じょうし (jōshi)}]{\ja{じょうし} (jōshi) — a boss; one's superior}

\label{lesson:JA-C230-joushi}



\begin{quote}
Before the new one: say the Japanese for a colleague, then the Japanese for a company president.
\end{quote}

\subsection*{You'll want to know: \ja{じょうし}}
\textbf{\ja{じょうし}} — \emph{jōshi} — \textquotedblleft{}a boss; one's superior\textquotedblright{}.

\begin{quote}
\textbf{\ja{じょうしは} \ja{いつも} \ja{ていねいに} \ja{おしえます。}} — \emph{jōshi wa itsumo teinei ni oshiemasu} — \textquotedblleft{}My boss always explains things carefully.\textquotedblright{}
\end{quote}

\begin{cousinweb}
\textbf{\ja{じょうし}} (\emph{jōshi}) is Sino-Japanese (kango): two Chinese characters read with their on'yomi, readings borrowed from Chinese. They are \emph{jō}, \textquotedblleft{}above\textquotedblright{} (Middle Chinese \emph{dzyangH}) and \emph{shi}, \textquotedblleft{}official\textquotedblright{} (\emph{si}) (KANJIDIC2; Middle Chinese in Baxter's transcription).
\end{cousinweb}

\begin{grammarlens}[title={What you've built}]
One more word for everyday life at A2.
\end{grammarlens}

\subsection*{Your turn}
\begin{itemize}
\raggedright
  \item \emph{Say it:} \emph{jōshi}
  \item \emph{Say it:} \emph{jōshi}, once more
  \item \emph{Say it:} \emph{jōshi wa itsumo teinei ni oshiemasu} — My boss always explains things carefully.
  \item \emph{Recall:} say the Japanese for a colleague, then the Japanese for a company president, then say \emph{jōshi} again
\end{itemize}

\begin{tcolorbox}[breakable,colback=teal!4,colframe=teal!35!black,title={Before you move on}]
What does \ja{じょうし} mean? (\textquotedblleft{}A boss; one's superior\textquotedblright{}.) Say it once more.
\end{tcolorbox}
\section[\texorpdfstring{\ja{せんぱい} (senpai)}{せんぱい (senpai)}]{\ja{せんぱい} (senpai) — a senior; an older schoolmate or colleague}

\label{lesson:JA-C230-senpai}



\begin{quote}
Before the new one: say the Japanese for a company president, then the Japanese for a boss.
\end{quote}

\subsection*{You'll want to know: \ja{せんぱい}}
\textbf{\ja{せんぱい}} — \emph{senpai} — \textquotedblleft{}a senior; an older schoolmate or colleague\textquotedblright{}.

\begin{quote}
\textbf{\ja{せんぱいは} \ja{よく} \ja{しゅくだいを} \ja{てつだいます。}} — \emph{senpai wa yoku shukudai o tetsudaimasu} — \textquotedblleft{}My senior often helps with my homework.\textquotedblright{}
\end{quote}

\begin{cousinweb}
\textbf{\ja{せんぱい}} (\emph{senpai}) is Sino-Japanese (kango): two Chinese characters read with their on'yomi, readings borrowed from Chinese. They are \emph{sen}, \textquotedblleft{}before\textquotedblright{} (Middle Chinese \emph{sen}) and \emph{pai}, \textquotedblleft{}fellows\textquotedblright{} (\emph{pwojH}) (KANJIDIC2; Middle Chinese in Baxter's transcription).
\end{cousinweb}

\begin{grammarlens}[title={What you've built}]
One more word for everyday life at A2.
\end{grammarlens}

\subsection*{Your turn}
\begin{itemize}
\raggedright
  \item \emph{Say it:} \emph{senpai}
  \item \emph{Say it:} \emph{senpai}, once more
  \item \emph{Say it:} \emph{senpai wa yoku shukudai o tetsudaimasu} — My senior often helps with my homework.
  \item \emph{Recall:} say the Japanese for a company president, then the Japanese for a boss, then say \emph{senpai} again
\end{itemize}

\begin{tcolorbox}[breakable,colback=teal!4,colframe=teal!35!black,title={Before you move on}]
What does \ja{せんぱい} mean? (\textquotedblleft{}A senior; an older schoolmate or colleague\textquotedblright{}.) Say it once more.
\end{tcolorbox}
\section[\texorpdfstring{\ja{こうはい} (kōhai)}{こうはい (kōhai)}]{\ja{こうはい} (kōhai) — a junior; a younger schoolmate or colleague}

\label{lesson:JA-C230-kouhai}



\begin{quote}
Before the new one: say the Japanese for a boss, then the Japanese for a senior.
\end{quote}

\subsection*{You'll want to know: \ja{こうはい}}
\textbf{\ja{こうはい}} — \emph{kōhai} — \textquotedblleft{}a junior; a younger schoolmate or colleague\textquotedblright{}.

\begin{quote}
\textbf{\ja{こうはいは} \ja{よく} \ja{せんぱいに} \ja{しつもんを} \ja{します。}} — \emph{kōhai wa yoku senpai ni shitsumon o shimasu} — \textquotedblleft{}The juniors often ask their seniors questions.\textquotedblright{}
\end{quote}

\begin{cousinweb}
\textbf{\ja{こうはい}} (\emph{kōhai}) is Sino-Japanese (kango): two Chinese characters read with their on'yomi, readings borrowed from Chinese. They are \emph{kō}, \textquotedblleft{}after, behind\textquotedblright{} (Middle Chinese \emph{huwX}) and \emph{hai}, \textquotedblleft{}fellows\textquotedblright{} (\emph{pwojH}) (KANJIDIC2; Middle Chinese in Baxter's transcription).
\end{cousinweb}

\begin{grammarlens}[title={What you've built}]
One more word for everyday life at A2.
\end{grammarlens}

\subsection*{Your turn}
\begin{itemize}
\raggedright
  \item \emph{Say it:} \emph{kōhai}
  \item \emph{Say it:} \emph{kōhai}, once more
  \item \emph{Say it:} \emph{kōhai wa yoku senpai ni shitsumon o shimasu} — The juniors often ask their seniors questions.
  \item \emph{Recall:} say the Japanese for a boss, then the Japanese for a senior, then say \emph{kōhai} again
\end{itemize}

\begin{tcolorbox}[breakable,colback=teal!4,colframe=teal!35!black,title={Before you move on}]
What does \ja{こうはい} mean? (\textquotedblleft{}A junior; a younger schoolmate or colleague\textquotedblright{}.) Say it once more.
\end{tcolorbox}
\section[\texorpdfstring{\ja{がんばる} (ganbaru)}{がんばる (ganbaru)}]{\ja{がんばる} (ganbaru) — to do one's best; to keep at it}

\label{lesson:JA-C230-ganbaru}



\begin{quote}
Before the new one: say the Japanese for a senior, then the Japanese for a junior.
\end{quote}

\subsection*{You'll want to know: \ja{がんばる}}
\textbf{\ja{がんばる}} — \emph{ganbaru} — \textquotedblleft{}to do one's best; to keep at it\textquotedblright{}.

\begin{quote}
\textbf{\ja{こうはいは} \ja{まいにち} \ja{しごとを} \ja{がんばります。}} — \emph{kōhai wa mainichi shigoto o ganbarimasu} — \textquotedblleft{}The junior does his best at work every day.\textquotedblright{}
\end{quote}

Polite \textbf{\ja{がんばります}}, \emph{ganbarimasu}; a godan verb.

\begin{cousinweb}
\textbf{\ja{がんばる}} (\emph{ganbaru}) is written with two characters meaning \textquotedblleft{}stubborn\textquotedblright{} and \textquotedblleft{}stretch\textquotedblright{}, read as one verb in \emph{-ru}; its older origin is not settled (JMdict, KANJIDIC2).
\end{cousinweb}

\begin{grammarlens}[title={What you've built}]
One more word for everyday life at A2.
\end{grammarlens}

\subsection*{Your turn}
\begin{itemize}
\raggedright
  \item \emph{Say it:} \emph{ganbaru}
  \item \emph{Say it:} \emph{ganbaru}, once more
  \item \emph{Say it:} \emph{kōhai wa mainichi shigoto o ganbarimasu} — The junior does his best at work every day.
  \item \emph{Recall:} say the Japanese for a senior, then the Japanese for a junior, then say \emph{ganbaru} again
\end{itemize}

\begin{tcolorbox}[breakable,colback=teal!4,colframe=teal!35!black,title={Before you move on}]
What does \ja{がんばる} mean? (\textquotedblleft{}To do one's best; to keep at it\textquotedblright{}.) Say it once more.
\end{tcolorbox}
